% Adapted from an OpenAI Math manuscript. See PROVENANCE.md.
\section{Arbitrary convex bodies and consequences}
\label{sec:approximation}

The passage from finite strip bodies to arbitrary convex bodies uses an approximation in the polar. This keeps the momentum factor inside the desired polar and enlarges the position factor by an arbitrarily small amount.

\begin{lemma}
\label{lem:polar-approximation}
Let $K\subset\mathbb{R}^n$ be an origin-symmetric convex body. For every $\alpha>0$, there are finitely many vectors $b_1,\ldots,b_m\in\mathbb{R}^n$ spanning $\mathbb{R}^n$ such that
\[
 K'=\{x\in\mathbb{R}^n:|\pairing{b_j}{x}|\le1,\quad 1\le j\le m\}
\]
satisfies
\begin{equation}\label{eq:polar-approximation}
 K\subset K'\subset(1+\alpha)K,
 \qquad (K')^\circ\subset K^\circ.
\end{equation}
\end{lemma}

\begin{proof}
Put $\lambda=1+\alpha$. The support function of $K^\circ$ has a positive minimum $\rho$ on the unit sphere. Choose a finite $\delta$-net $\{b_j\}$ in $K^\circ$, with $0<\delta\leq(1-\lambda^{-1})\rho$, and let $P=\operatorname{conv}\{\pm b_j\}$. At each unit vector $v$, choose a maximizer for $h_{K^\circ}(v)$ and approximate it by one of the $b_j$. This gives
\[
 h_P(v)\geq h_{K^\circ}(v)-\delta\geq\lambda^{-1}h_{K^\circ}(v).
\]
The support-function characterization of inclusion gives $\lambda^{-1}K^\circ\subset P\subset K^\circ$. In particular the $b_j$ span the ambient space. Taking polars reverses inclusions, and $K'=P^\circ$ has the stated strip description.
\end{proof}

\begin{proof}[Proof of \cref{thm:main}]
Fix $0<c<4$. Choose $\alpha>0$ such that $\lambda c<4$, where $\lambda=1+\alpha$, and choose one body $K'$ from \cref{lem:polar-approximation}. \Cref{prop:finite-strips} gives a smooth symplectic embedding
\[
 e:B^{2n}(\lambda c)\longrightarrow
 \operatorname{int} K'\times\operatorname{int}\bigl((K')^\circ\bigr).
\]
If one body is contained in another, every interior ball in the first is contained in the second. Thus \eqref{eq:polar-approximation} implies
\[
 \operatorname{int} K'\subset\lambda\operatorname{int} K,
 \qquad \operatorname{int}\bigl((K')^\circ\bigr)\subset\operatorname{int} K^\circ,
\]
even when their boundaries meet. Regard $e$ as an embedding into $\lambda\operatorname{int} K\times\operatorname{int} K^\circ$. Define
\[
 D_{\sqrt\lambda}(z)=\sqrt\lambda\,z,
 \qquad Q_\lambda(q,p)=(q/\lambda,p).
\]
The composition
\begin{equation}\label{eq:general-rescaling}
 Q_\lambda\circ e\circ D_{\sqrt\lambda}:
 B^{2n}(c)\longrightarrow U_K
\end{equation}
is a smooth embedding. Since
\[
 D_{\sqrt\lambda}^*\omega_0=\lambda\omega_0,
 \qquad Q_\lambda^*\omega_0=\lambda^{-1}\omega_0,
\]
its pullback of $\omega_0$ equals $\omega_0$. This proves $c_G(U_K)\ge4$.

For the upper bound, we use the supporting-slab construction of Artstein-Avidan, Karasev, and Ostrover \cite[Remark~4.2]{AAKO2014}. Choose $q_0\in K$ and $p_0\in K^\circ$ with $\pairing{q_0}{p_0}=1$. Such a pair is obtained by maximizing any nonzero linear functional on $K$ and normalizing its maximum to $1$. Complete $q_0$ to a basis of $\mathbb{R}^n$ by vectors in $\ker p_0$, and let $A$ be the matrix of this basis. The linear change
\[
 Q=A^{-1}q,\qquad P=A^Tp
\]
is symplectic and has first coordinates $Q_1=\pairing{p_0}{q}$ and $P_1=\pairing{q_0}{p}$. Both functionals are nonzero, so symmetry, polarity, and the interior condition give $|Q_1|<1$ and $|P_1|<1$ on $U_K$. Thus this change embeds $U_K$ into $(-1,1)^2\times\mathbb{R}^{2n-2}$, with the square in its first conjugate coordinate pair.

For completeness, this open square is area-preservingly diffeomorphic to the disk of area $4$. Put $R=2/\sqrt\pi$, $h(X)=\sqrt{R^2-X^2}$ for $-R<X<R$, and
\[
 H(X)=-1+\int_{-R}^{X}h(s)\,\rd s.
\]
Since $\int_{-R}^{R}h(s)\,\rd s=2$ and $H'=h>0$, $H$ is a smooth diffeomorphism from $(-R,R)$ onto $(-1,1)$. The map
\[
 (u,v)\longmapsto
 \bigl(X=H^{-1}(u),\ Y=v\,h(H^{-1}(u))\bigr)
\]
has image $B^2(4)$ and Jacobian $X'(u)h(X)=1$. Applying it to the first conjugate pair gives a symplectic embedding $U_K\hookrightarrow B^2(4)\times\mathbb{R}^{2n-2}$. Gromov's nonsqueezing theorem \cite{Gromov1985} therefore bounds the capacity of every ball in $U_K$ by $4$. Combined with the lower bound, this proves $c_G(U_K)=4$.
\end{proof}

\begin{remark}
\label{rem:parameter-order}
For each prescribed $0<c<4$, first choose $\alpha>0$ so that $\lambda c<4$, where $\lambda=1+\alpha$. Then fix the finite body $K'$ and its $m$ defining constraints. Choose $\eta>0$ so that $(1+\eta)\lambda c<4$, then choose $k$ large enough that
\[
 \frac{C_{K'}}k+m\epsilon_k<\frac{\eta\pi}{4}.
\]
This is possible because the constraint list is fixed before $k$ varies. \Cref{prop:finite-strips}, applied at capacity $\lambda c$, and the final rescaling then give the required ball. Each prescribed capacity is obtained by one finite construction; no limit of embeddings is required.
\end{remark}

\begin{corollary}
\label{cor:mahler}
For every integer $n\ge1$ and every origin-symmetric convex body $K\subset\mathbb{R}^n$,
\[
 \vol_n(K)\,\vol_n(K^\circ)\ge\frac{4^n}{n!}.
\]
\end{corollary}

\begin{proof}
For $n=1$, write $K=[-a,a]$ with $a>0$. Then $K^\circ=[-a^{-1},a^{-1}]$, and the product of their lengths is $4$. Assume henceforth that $n\ge2$.

We first justify that taking interiors does not change the volumes in the statement. If a convex body $C\subset\mathbb{R}^n$ contains $B(0,r)$, then convexity gives, for $0<t<1$,
\[
 tC+B(0,(1-t)r)\subset C.
\]
Hence $tC\subset\operatorname{int} C$, and
\[
 t^n\vol_n(C)\le\vol_n(\operatorname{int} C)\le\vol_n(C).
\]
Letting $t\uparrow1$ proves $\vol_n(\operatorname{int} C)=\vol_n(C)$. Apply this to $K$ and $K^\circ$, both of which contain a neighborhood of the origin. The product formula for Lebesgue measure then gives
\begin{equation}\label{eq:polar-product-volume}
 \vol_{2n}(U_K)=\vol_n(K)\,\vol_n(K^\circ).
\end{equation}

For every $0<c<4$, \cref{thm:main} supplies a smooth symplectic embedding $f:B^{2n}(c)\to U_K$. Taking the $n$th exterior power of $f^*\omega_0=\omega_0$ shows that $\det Df=1$, so the change of variables formula gives
\[
 \frac{c^n}{n!}
 =\vol_{2n}\bigl(B^{2n}(c)\bigr)
 =\vol_{2n}\bigl(f(B^{2n}(c))\bigr)
 \le\vol_{2n}(U_K).
\]
Here the first equality follows from the Euclidean ball volume $\pi^nR^{2n}/n!$ at radius $R=\sqrt{c/\pi}$. Combining this with \eqref{eq:polar-product-volume} and letting $c\uparrow4$ proves the claimed inequality.
\end{proof}

\begin{remark}\label{rem:cube-volume-product}
The cube $[-1,1]^n$ has volume $2^n$. Its polar is $\{p:\sum_{j=1}^n|p_j|\le1\}$, whose intersection with each coordinate orthant is a simplex of volume $1/n!$. These $2^n$ simplices meet only along faces, so the polar has volume $2^n/n!$. Their volume product is therefore $4^n/n!$, showing that the constant in \cref{cor:mahler} is sharp.
\end{remark}

\begin{corollary}
\label{cor:polar-packings}
Let $n\ge3$ and $k\ge1$ be integers, and let $K\subset\mathbb{R}^n$ be an origin-symmetric convex body. Suppose that $R_1,\ldots,R_k>0$ satisfy
\[
 \sum_{i=1}^k R_i^n<4^n,
 \qquad R_i+R_j<4\quad(1\le i<j\le k).
\]
Then the closed balls $\overline{B^{2n}(R_i)}$ admit symplectic embeddings into $U_K$ with pairwise disjoint images, each defined on an open neighborhood of the corresponding closed ball.
\end{corollary}

\begin{proof}
Strictness and finiteness permit a choice of $0<c<4$ such that $\sum_i R_i^n<c^n$ and $R_i+R_j<c$ for every $i<j$. The ball-packing theorem \cite[Theorem~1.1]{OpenAIBallPackings2026} embeds these closed balls into $B^{2n}(c)$, with embeddings defined on neighborhoods. Their compact images are pairwise disjoint and lie in the open target, so the neighborhoods can be shrunk to have pairwise disjoint images contained in $B^{2n}(c)$. Composing with the embedding in \cref{thm:main} proves the assertion.
\end{proof}
